\begin{historia}[1]{HU-20}{Registrar una tarea de arrastrar y soltar}{Administrador}{1}{13}{\textbf{Como} administrador, \textbf{quiero} registrar una tarea de arrastrar y soltar, que consta de un título, un enunciado, la pregunta, una imagen de fondo y las piezas que el estudiante debe colocar en su lugar sobre ella, \textbf{para} contar en el banco de preguntas con tareas donde hay que mover objetos a su lugar.}
\hucriterio{Se registran el título de la tarea, su enunciado y la pregunta que se le hace al estudiante.}
\hucriterio{El enunciado admite texto con formato e imágenes de tamaño ajustable, incluso imágenes pequeñas dentro de una línea de texto, como los dibujos del cuadernillo.}
\hucriterio{Se registran la imagen de fondo y una imagen por cada pieza.}
\hucriterio{El lugar correcto de cada pieza se define soltándola sobre el fondo.}
\hucriterio{El tamaño del fondo y de cada pieza se puede ajustar, y el editor muestra el tamaño que verá el estudiante.}
\hucriterio{Es posible marcar lugares adicionales donde se puede soltar una pieza aunque no sea la respuesta.}
\hucriterio{Si la tarea tiene más de una solución, se pueden registrar otras respuestas correctas.}
\hucriterio{Dos piezas con la misma imagen se consideran iguales: es indistinto cuál ocupa cada lugar.}
\hucriterio{Antes de guardar, el sistema lista los datos obligatorios que faltan, como el título, la pregunta o la respuesta correcta, y la lista se actualiza a medida que se completan.}
\end{historia}
